\documentclass[12pt,a4paper]{article}
\usepackage[utf8]{inputenc}
\usepackage[T1]{fontenc}
\usepackage{geometry}
\usepackage{longtable}
\usepackage{array}
\usepackage{hyperref}
\geometry{a4paper, margin=2.5cm}
\title{\textbf{Plano de Testes}\\Sistema Reserva de Salas de Estudo}
\author{Ailton Baren Pereira da Silva}
\date{Junho de 2026}
\begin{document}
\maketitle

\section{Objetivo}
Validar se o sistema atende aos principais fluxos funcionais, regras de negócio e restrições de acesso. O plano contempla testes automatizados com \texttt{pytest} e validação manual para demonstração acadêmica.

\section{Ambiente}
\begin{itemize}
    \item Python 3, Flask, Flask-SQLAlchemy e Flask-Login.
    \item PostgreSQL como banco da aplicação.
    \item Jinja2 e Bootstrap 5 por CDN.
    \item \texttt{pytest} para testes automatizados.
\end{itemize}

\section{Estratégia}
Os testes automatizados utilizam o cliente de testes do Flask e banco isolado em memória, permitindo validar as regras de negócio sem depender do banco PostgreSQL de desenvolvimento.

\section{Escopo}
\subsection{Incluído}
\begin{itemize}
    \item Login válido e inválido.
    \item Acesso protegido e bloqueio de área administrativa para aluno.
    \item Cadastro e edição de usuários pelo administrador.
    \item Consulta de salas e horários.
    \item Indicação de sala bloqueada no dia atual na lista de salas.
    \item Criação, rejeição e cancelamento de reservas.
    \item Cadastro, edição e exclusão controlada de salas.
    \item Bloqueio de salas por período.
    \item Configuração de limite de reservas ativas.
    \item Renderização das principais telas.
\end{itemize}

\subsection{Não incluído}
\begin{itemize}
    \item Testes de carga.
    \item Testes em múltiplos navegadores.
    \item Testes de segurança avançados.
    \item Testes com serviços externos.
\end{itemize}

\section{Casos de Teste}
\begin{longtable}{|p{1.6cm}|p{6.2cm}|p{5.4cm}|p{1.8cm}|}
\hline
\textbf{ID} & \textbf{Caso} & \textbf{Resultado esperado} & \textbf{Status} \\ \hline
CT-01 & Login com credenciais válidas & Usuário autenticado & Aprovado \\ \hline
CT-02 & Login com senha inválida & Mensagem de erro exibida & Aprovado \\ \hline
CT-03 & Página protegida sem login & Redirecionamento para login & Aprovado \\ \hline
CT-04 & Aluno acessa painel admin & Erro 403 & Aprovado \\ \hline
CT-05 & Admin cadastra usuário & Usuário salvo e apto a login & Aprovado \\ \hline
CT-06 & Admin cadastra e-mail duplicado & Cadastro rejeitado & Aprovado \\ \hline
CT-07 & Admin edita usuário sem senha nova & Senha anterior preservada & Aprovado \\ \hline
CT-08 & Criação de reserva válida & Reserva ativa criada & Aprovado \\ \hline
CT-09 & Reserva em data passada & Reserva rejeitada & Aprovado \\ \hline
CT-10 & Hora final anterior à inicial & Reserva rejeitada & Aprovado \\ \hline
CT-11 & Reserva com duração diferente de uma hora & Reserva rejeitada & Aprovado \\ \hline
CT-12 & Conflito de horário da sala & Reserva rejeitada & Aprovado \\ \hline
CT-13 & Cancelamento com motivo & Reserva cancelada & Aprovado \\ \hline
CT-14 & Cancelamento sem motivo & Cancelamento rejeitado & Aprovado \\ \hline
CT-15 & Horários com duração fixa de uma hora & Horários de uma hora exibidos & Aprovado \\ \hline
CT-16 & Horário reservado mostra finalidade & Finalidade exibida & Aprovado \\ \hline
CT-17 & Admin cadastra sala válida & Sala salva & Aprovado \\ \hline
CT-18 & Sala com capacidade inválida & Sala rejeitada & Aprovado \\ \hline
CT-19 & Bloqueio por período & Bloqueio ativo criado & Aprovado \\ \hline
CT-20 & Bloqueio impede reserva & Reserva rejeitada & Aprovado \\ \hline
CT-21 & Admin cancela bloqueio & Bloqueio inativo & Aprovado \\ \hline
CT-22 & Admin configura limite & Configuração salva & Aprovado \\ \hline
CT-23 & Aluno excede limite & Reserva rejeitada & Aprovado \\ \hline
CT-24 & Lista de salas com bloqueio do dia atual & Sala exibida como bloqueada hoje & Aprovado \\ \hline
CT-25 & Admin exclui sala sem histórico & Sala removida & Aprovado \\ \hline
CT-26 & Admin tenta excluir sala com reserva & Exclusão rejeitada & Aprovado \\ \hline
\end{longtable}

\section{Execução}
Comando utilizado:
\begin{verbatim}
source .venv/bin/activate
pytest
\end{verbatim}

\section{Resultado}
Última execução registrada:
\begin{verbatim}
36 passed
\end{verbatim}

\section{Critérios de Aceite}
O sistema é aprovado nos testes quando todos os testes automatizados passam, as regras de reserva são respeitadas, o acesso administrativo é protegido e as mensagens são exibidas em português.

\end{document}
